\begin{convention}[Index suppression and aggregation]
\label{conv:notation}

\paragraph{Root variables, projected indices, and suppressed indices.}
Every variable in the model is defined over a fixed tuple of indices, called
its \emph{root index set}; all variables listed in
Table~\ref{tab:nomenclature} are \emph{root variables}.

Given a root variable $\phi_{\alpha}$ with full index tuple
$\alpha = (\alpha_1, \dots, \alpha_k)$ ranging over
$\mathcal{I}_1 \times \cdots \times \mathcal{I}_k$, partition the positions
into a \emph{projected} subset and a \emph{suppressed} subset. Collecting
the projected positions into tuple $\pi$ and the suppressed positions into
tuple $\sigma$, so that $\alpha = (\pi, \sigma)$, defines two domains:
\begin{align*}
  \mathcal{P} &\;\coloneqq\;
    \prod_{\text{projected } i} \mathcal{I}_i
    & &\text{(projected domain, indexing the derived variable),}\\
  \mathcal{S} &\;\coloneqq\;
    \prod_{\text{suppressed } i} \mathcal{I}_i
    & &\text{(suppressed domain, collapsed by aggregation).}
\end{align*}
The root variable may then be written $\phi_{\pi,\sigma}$ to make both
parts of the index explicit. Writing $\phi$ with only $\pi$ denotes a
\emph{derived variable}, obtained by collapsing the \emph{group}
\[
  \mathcal{G}(\pi)
  \;\coloneqq\;
  \bigl\{\, \phi_{\pi,\sigma} : \sigma \in \mathcal{S} \bigr\},
  \qquad \forall\,\pi \in \mathcal{P},
\]
via an aggregation function chosen from Rules~1--3. The accent on $\phi$
identifies which rule is applied, and $|\mathcal{S}|$ denotes the
cardinality of the suppressed domain.

\paragraph{Dot notation for repeated index sets.}
When two or more indices range over the same set---as $i, j \in V$ do in
$NF_{cijmpr}$---suppressing one of them by position alone would be
ambiguous. A dot~$(\,\cdot\,)$ marks the suppressed position explicitly:
\[
  \operatorname{Agg}\!\bigl(\{NF_{cijmpr} : i \in V\}\bigr)
  \;\equiv\;
  \operatorname{Agg}\!\bigl(NF_{\cdot jmpr}\bigr),
  \qquad
  \operatorname{Agg}\!\bigl(\{NF_{cijmpr} : j \in V\}\bigr)
  \;\equiv\;
  \operatorname{Agg}\!\bigl(NF_{i\cdot mpr}\bigr).
\]
This convention applies uniformly to all three rules below.

%──────────────────────────────────────────────────────────────────────────────
\medskip
\noindent\textbf{Rule 1 (Sum aggregation,~$\widetilde{\phi}$).}

\smallskip
\noindent\textit{Abstract definition.}
Let $\phi_{\pi,\sigma}$ be any root variable (continuous, integer, or
binary) with projected domain $\mathcal{P}$ and suppressed domain
$\mathcal{S}$. The sum-derived variable is
\[
  \widetilde{\phi}_{\pi}
  \;\coloneqq\;
  \sum_{\sigma \in \mathcal{S}} \phi_{\pi,\sigma},
  \qquad \forall\,\pi \in \mathcal{P}.
\]
For binary root variables the result is an \emph{integer count} of active
members in the group; it is not itself binary. When $\widetilde{\phi}_{\pi}$
is introduced explicitly as a model variable, its definition is enforced by
the equality constraint
\[
  \widetilde{\phi}_{\pi} = \sum_{\sigma \in \mathcal{S}} \phi_{\pi,\sigma},
  \qquad \forall\,\pi \in \mathcal{P}.
\]

\smallskip
\noindent\textit{Example.}
Consider the root variable $NF_{cijmpr}$ with full index tuple
$(c,i,j,m,p,r) \in C \times V \times V \times \mathcal{M} \times P \times R$.

\begin{itemize}
  \item \textbf{Single suppression.}
        Projected tuple $\pi = (c,i,j,p,r)$, suppressed tuple
        $\sigma = m$, suppressed domain $\mathcal{S} = \mathcal{M}$,
        projected domain $\mathcal{P} = C \times V \times V \times P \times R$:
        \[
          \widetilde{NF}_{cijpr}
          \;\coloneqq\;
          \sum_{m \in \mathcal{M}} NF_{cijmpr},
          \qquad
          \forall\,(c,i,j,p,r) \in \mathcal{P}.
        \]

  \item \textbf{Multiple suppressions.}
        Projected tuple $\pi = (c,i,j)$, suppressed tuple
        $\sigma = (m,p,r)$, suppressed domain
        $\mathcal{S} = \mathcal{M} \times P \times R$,
        projected domain $\mathcal{P} = C \times V \times V$:
        \[
          \widetilde{NF}_{cij}
          \;\coloneqq\;
          \sum_{(m,p,r)\,\in\,\mathcal{S}}
          NF_{cijmpr},
          \qquad
          \forall\,(c,i,j) \in \mathcal{P}.
        \]

  \item \textbf{Full suppression.}
        Projected tuple $\pi = \emptyset$ (scalar result), suppressed domain
        $\mathcal{S} = C \times V \times V \times \mathcal{M} \times P \times R$:
        \[
          \widetilde{NF}
          \;\coloneqq\;
          \sum_{(c,i,j,m,p,r)\,\in\,\mathcal{S}}
          NF_{cijmpr}.
        \]
\end{itemize}

For the binary variable $X_{cijm}$, with projected tuple $\pi=(c,i,j)$,
suppressed tuple $\sigma = m$, and suppressed domain $\mathcal{S}=\mathcal{M}$:
\[
  \widetilde{X}_{cij}
  \;\coloneqq\;
  \sum_{m \in \mathcal{M}} X_{cijm}
  \;\in\; \{0, 1, \dots, |\mathcal{M}|\},
  \qquad \forall\,(c,i,j) \in C \times V \times V,
\]
counting how many modes are active for class~$c$ on lane $(i,j)$.

%──────────────────────────────────────────────────────────────────────────────
\medskip
\noindent\textbf{Rule 2 (OR aggregation,~$\widecheck{\phi}$).}

\smallskip
\noindent\textit{Abstract definition.}
Let $\phi_{\pi,\sigma}$ be a \emph{binary} root variable with projected
domain $\mathcal{P}$ and suppressed domain $\mathcal{S}$. For each
$\pi \in \mathcal{P}$, a new binary decision variable $\widecheck{\phi}_{\pi}$
is introduced and linked to the corresponding sum $\widetilde{\phi}_{\pi}$
via:
\begin{align}
  \widecheck{\phi}_{\pi} &\leq \widetilde{\phi}_{\pi},
    \label{eq:or-lb}\\
  \widetilde{\phi}_{\pi} &\leq |\mathcal{S}|\cdot\widecheck{\phi}_{\pi},
    \label{eq:or-ub}
\end{align}
for all $\pi \in \mathcal{P}$. Together,
\eqref{eq:or-lb}--\eqref{eq:or-ub} enforce
$\widecheck{\phi}_{\pi} = 1 \iff \widetilde{\phi}_{\pi} \geq 1$, i.e.\
at least one root binary in the group $\mathcal{G}(\pi)$ equals~1.

\smallskip
\noindent\textit{Example.}
For the binary variable $X_{cijm}$, with projected tuple $\pi = (c,i,j)$,
projected domain $\mathcal{P} = C \times V \times V$, suppressed tuple
$\sigma = m$, and suppressed domain $\mathcal{S} = \mathcal{M}$
($|\mathcal{S}| = |\mathcal{M}|$), the OR-derived binary
$\widecheck{X}_{cij}$ indicates whether \emph{any} mode is active for
class~$c$ on lane $(i,j)$:
\begin{align*}
  \widecheck{X}_{cij} &\leq \widetilde{X}_{cij}, \\
  \widetilde{X}_{cij} &\leq |\mathcal{M}|\cdot\widecheck{X}_{cij},
  \qquad \forall\,(c,i,j) \in \mathcal{P},
\end{align*}
so $\widecheck{X}_{cij} = 1$ if and only if at least one mode $m \in \mathcal{M}$
has $X_{cijm} = 1$.

%──────────────────────────────────────────────────────────────────────────────
\medskip
\noindent\textbf{Rule 3 (At-most-one aggregation,~$\widehat{\phi}$).}

\smallskip
\noindent\textit{Abstract definition.}
Let $\phi_{\pi,\sigma}$ be a \emph{binary} root variable with projected
domain $\mathcal{P}$ and suppressed domain $\mathcal{S}$. For each
$\pi \in \mathcal{P}$, a new binary decision variable $\widehat{\phi}_{\pi}$
is introduced and linked to the corresponding sum $\widetilde{\phi}_{\pi}$
via:
\begin{align}
  \widehat{\phi}_{\pi} &\leq \widetilde{\phi}_{\pi},
    \label{eq:amo-lb}\\
  \widetilde{\phi}_{\pi} &\leq \widehat{\phi}_{\pi},
    \label{eq:amo-ub}
\end{align}
for all $\pi \in \mathcal{P}$. Because $\widehat{\phi}_{\pi}$ is binary,
\eqref{eq:amo-lb}--\eqref{eq:amo-ub} force $\widetilde{\phi}_{\pi} \in \{0,1\}$,
asserting that \emph{at most one} root binary in the group $\mathcal{G}(\pi)$
may equal~1. Equivalently: $\widehat{\phi}_{\pi} = 1$ if and only if exactly
one root binary is active; $\widehat{\phi}_{\pi} = 0$ if and only if none are.

\smallskip
\noindent\textit{Example.}
For the binary variable $X_{cijm}$, with projected tuple $\pi = (c,i,j)$,
projected domain $\mathcal{P} = C \times V \times V$, suppressed tuple
$\sigma = m$, and suppressed domain $\mathcal{S} = \mathcal{M}$, the
at-most-one-derived binary $\widehat{X}_{cij}$ asserts that at most one
mode may be simultaneously active for class~$c$ on lane $(i,j)$:
\begin{align*}
  \widehat{X}_{cij} &\leq \widetilde{X}_{cij}, \\
  \widetilde{X}_{cij} &\leq \widehat{X}_{cij},
  \qquad \forall\,(c,i,j) \in \mathcal{P},
\end{align*}
so $\widehat{X}_{cij} = \widetilde{X}_{cij} \in \{0,1\}$, enforcing mutual
exclusivity across all modes $m \in \mathcal{M}$.

\begin{figure}[H]
\centering
\resizebox{0.95\linewidth}{!}{%
  % --- ADJUSTABLE PARAMETERS ---
\def\hsep{4.0}
\def\vsep{1.4}
\def\logicXshift{7.5}
% -----------------------------
\begin{tikzpicture}[
  every node/.style = {
    draw, rounded corners=3pt, inner sep=5pt,
    align=center, font=\small, thick
  },
  rootBox/.style   = {fill=blue!10, draw=blue!60!black},
  sumBox/.style    = {fill=orange!10, draw=orange!60!black},
  binBox/.style    = {fill=red!5, draw=red!80!black, dashed},
  arr/.style       = {->, >=stealth, thick},
  barr/.style      = {->, >=stealth, thick, dashed, red!80!black},
  lbl/.style       = {draw=none, fill=none, font=\footnotesize\itshape, inner sep=2pt},
  blbl/.style      = {draw=none, fill=none, font=\footnotesize\bfseries, inner sep=2pt, text=red!80!black},
  graylbl/.style   = {draw=none, fill=none, font=\small\bfseries, gray!50}
]

%% --- SUM AGGREGATION SECTION (Left) ---
\node[rootBox] (root) at (0,0) {$\phi_{cijm}$};
\node[sumBox] (Fcij)  at (-\hsep/2, \vsep)
  {$\widetilde{\phi}_{cij} \coloneqq \sum_{m \in \mathcal{M}} \phi_{cijm}$};
\node[sumBox] (Fijm)  at (\hsep/2, \vsep)
  {$\widetilde{\phi}_{ijm} \coloneqq \sum_{c \in C} \phi_{cijm}$};
\node[sumBox] (Fij)   at (0, 2*\vsep)
  {$\widetilde{\phi}_{ij} \coloneqq \sum_{c \in C} \sum_{m \in \mathcal{M}} \phi_{cijm}$};

\draw[arr] (root.north) -- (Fcij.south);
\draw[arr] (root.north) -- (Fijm.south);
\draw[arr] (Fcij.north) -- (Fij.south);
\draw[arr] (Fijm.north) -- (Fij.south);

% Dot notation
\node[sumBox] (dotj) at (-\hsep/2, -\vsep)
  {$\widetilde{\phi}_{c \cdot j m} \coloneqq \sum_{i \in V} \phi_{cijm}$};
\node[sumBox] (doti) at (\hsep/2, -\vsep)
  {$\widetilde{\phi}_{ci \cdot m} \coloneqq \sum_{j \in V} \phi_{cijm}$};

\draw[arr] (root.south) -- (dotj.north);
\draw[arr] (root.south) -- (doti.north);

%% --- BINARY LOGIC AGGREGATION SECTION (Right) ---
\begin{scope}[xshift=\logicXshift cm]
\node[rootBox] (Xroot) at (0, 0) {$X_{cijm}$};
\node[binBox] (XOR)  at (0, \vsep*1.36)
  {$\widecheck{X}_{cij}$ (OR aggregate) \\
   \scriptsize $\widecheck{X}_{cij} \leq \widetilde{X}_{cij}$ \\
   \scriptsize $\widetilde{X}_{cij} \leq |\mathcal{M}|\,\widecheck{X}_{cij}$};
\node[binBox] (XHAT) at (0, -\vsep*1.36)
  {$\widehat{X}_{cij}$ (At-most-one) \\
   \scriptsize $\widehat{X}_{cij} \leq \widetilde{X}_{cij}$ \\
   \scriptsize $\widetilde{X}_{cij} \leq \widehat{X}_{cij}$};

\draw[barr] (Xroot.north) -- (XOR.south)  node[blbl, midway, right] {Rule 3};
\draw[barr] (Xroot.south) -- (XHAT.north) node[blbl, midway, right] {Rule 4};
\end{scope}

%% --- DIVIDER AND LABELS ---
\draw[thick, gray!30, dashed]
  (\logicXshift/2, -1.5*\vsep) -- (\logicXshift/2, 2.8*\vsep);
\node[graylbl] at (0,            2.7*\vsep) {Sum Aggregation};
\node[graylbl] at (\logicXshift, 2.7*\vsep) {Binary Logic Aggregation};

\end{tikzpicture}
}
\caption{Derivation lattice: sum aggregation (left) and binary logic
  aggregation (right).}
\label{fig:index-lattice}
\end{figure}
\end{convention}
